\section{淝水前后：北方重组与东晋的终局}
\label{sec:eastern-jin-fei-river-liu-yu}

四世纪中叶以后的北方并非只是建康朝廷北伐的背景。后赵瓦解后，冉闵在邺城称魏，内战中针对羯胡及后赵旧部的暴力清洗留下极其沉重的记录；群体名称、参与范围和伤亡数字都须谨慎，不能把史书分类变成固定的族群本质。与此同时，苻健在长安建立前秦，慕容氏在东北和中原扩展前燕。多个政权同时重组人口、军队和城郭，洛阳、长安、邺城与河西的归属因而不断改变。

\subsection{前秦的集中与东晋的北望}

苻坚三五七年即位，任用王猛整顿吏治、抑制豪强；政策条文和“集权”的程度不能由后世称颂直接外推到每一郡县，但前秦确实逐步获得在关中征发、任官和调兵的能力。三七〇年王猛攻灭前燕，三七三年又取东晋梁、益二州北部州郡；三七六年兼并前凉、灭代国。每一次“灭国”之后的地方控制速度、部众去向和官署重建都不相同，不能用连续的胜利名称抹平差别；不过，前秦掌握的资源和道路已足以把东晋置于汉水、淮河与四川北缘的多重压力之下。

东晋并非没有北方目标。桓温数次北伐，三五六年进入洛阳并修复晋帝陵，三六九年又在枋头失利；收复旧都的象征不能取代稳定的补给和驻军。三七八年谢玄在京口训练北府兵，既利用侨民与北来军人的军事经验，也把淮南防务重新组织为可调动的力量。北府兵的实际构成和规模不宜由一个名称推成固定的社会集团，但它说明建康必须在迁徙、地方军镇和中央任命之间重新寻找兵源。

\subsection{淝水不是神话的捷径}

三八二年苻坚决定南伐，调集北方多地兵员与粮运；\eventanchor{JH-0383-FEI-RIVER}{东晋·淝水之战}次年谢玄、谢石等东晋军在淝水击败前秦军，苻坚北逃。战役结果可靠，所谓草木皆兵、风声鹤唳以及双方兵数、临阵口令和撤退经过，经过后世叙事反复雕琢，不能写成可逐帧复原的现场。较稳的解释不在于把胜负归给一条妙计，而在于前秦的远距离调兵、不同部众间的协调和淮河前线的补给都极为脆弱；东晋则在近于本土的水陆节点上，保有较能协同的军队。

淝水之后，前秦没有立即被一个胜者接收。三八四年慕容垂在河北起兵建后燕，姚苌在北地称秦；次年苻坚被杀。政权解体的速度、各部何时倒向和战斗细节在《晋书》〈载记〉与十六国史辑佚材料中屡有异文，却足以表明，大规模征服没有消除地方军政网络，反而使它们在中枢受挫时迅速取得独立空间。

\subsection{北魏的城郭与东晋的最后军镇}

三八六年拓跋珪重建代国，旋改国号魏；三九五年参合陂取胜，三九六年称帝，三九七年攻取中山，次年迁都平城。部落、城郭、汉官与北方农耕人口之间的关系并不能用“汉化”或“鲜卑化”一个词概括。迁都平城首先显示，北魏正在把代北的军事网络接入一个可以设官、储粮和发令的都城空间；它并不证明北方已被同一种制度均匀覆盖。

东晋内部也在失去这种接合能力。三九九年孙恩在会稽起兵，沿海州郡反复征战；四〇二年桓玄自荆州东下控制建康，次年废安帝自立，四〇四年刘裕等人恢复晋室。人物彼此的忠诚和叛变常被后出的传记写得过于整齐，城市、江道和外镇的控制权才是这些政变能够发生的条件。刘裕在四〇九年至四一三年间平南燕、收荆州、灭谯蜀和卢循余部，逐渐把长江中游和部分地方军镇纳入自己的军政网络。

四一六年刘裕北伐，次年入洛阳、长安，后秦灭亡；四一八年赫连勃勃乘其南返取长安，东晋的关中据点很快失去。前线为何不能久守，不能只用一位统帅南返或一支军队的道德来解释：河南、关中、江南之间的补给、地方接应与军镇利益都在限制选择。四二〇年刘裕受禅建宋，东晋至此结束。南方的皇统换了名称，北方的多政权竞争却仍在继续；两者共同构成此后南北朝的起点。

\subsection{史料的交叉与沉默}

《晋书》帝纪、列传和〈载记〉提供前秦、东晋、后燕与北魏兴起的主线，《资治通鉴》有助于校看并行年序；二者都以较晚的编纂和晋朝视野组织北方材料。十六国史辑佚可补国号、地名与互动线索，却常脱离原本上下文。王兴之夫妇墓志、爨宝子碑、炳灵寺题记以及城址与墓葬，为官职、地方网络和宗教活动留下不同于正史的锚点，仍不能替代战争、人口与行政实施的全景。本节据此保留政权转换和道路、军镇、都城的关系，将族群分类、兵额和地方控制范围留在可争论的材料边界内。

% 深描扩写：资源、道路与人的选择
\section{在事件之间：资源、道路与人的选择}
\subsection{苻坚与王猛的关中}
从一件可见之物进入，才不会把这段历史写成抽象的王朝更替。 三五七年苻坚即位，王猛参与整顿。长安、关中平原和河东陇右道路中的簿册、仓粟、城门和任命不是背景，而是命令、粮食、人员和消息能够移动的条件。苻坚、王猛、官员、豪强、军人、编户被放在同一条不平坦的关系链上：有人借它发布号令，有人依它转运供给，也有人在它被截断时改变依附。政策条文不等于全境执行，关中征发能力的增强确实改变扩张条件。

若只把三五七年苻坚即位，王猛参与整顿归结为一位皇帝、将领或谋臣的性格，便会错过更长的压力。继承安排、军镇分布、土地关系、边防开支和迁徙，往往早已改变了地方的承受能力；眼前的冲突只是把它们集中显露。有人能守城，有人能运粮，也有人只剩下离开原处的道路。 因而名义、武力和供给必须并置叙述，不能把其中任何一项当成万能的解释。

从地方尺度看，事情并不只发生在帝纪最常书写的宫城或战场。守令、吏员、舟人、农户、工匠、商人和随军家属未必留下姓名，却要承担征发、迁移、停市、守城或重登簿册的后果。仓储是否有粮，渡口是否可用，旧籍能否辨认，新来者与旧居民能否取得暂时位置，决定命令能否成为日常秩序。

人事转折也应在这种约束中理解。苻坚、王猛、官员、豪强、军人、编户不是脱离制度而自由行动的人物，他们受官号、部众、家族、道路、既得任命以及对对手力量的有限消息所限。传记常把胜者写得果断、败者写得必然失败；它保存冲突轮廓，却容易遮住联盟为何在此刻成立，又为何在城防、补给或继承改变后瓦解。

三五七年苻坚即位，王猛参与整顿的直接结果，通常不是把一片地方从混乱变为整齐。它会重写谁守备、哪些户口重登、哪条水道由谁巡护、哪些旧官仍可使用，也会让一部分人带着家属、技术、经卷或军籍离开原位。新获土地需要守军，新接军队需要粮饷，新立皇统需要文书；这些持续成本很快进入下一场危机，形成难以回避的反馈回路。

这也是制度得失必须落在具体人事上的原因。能够提高国家可见性的制度，可能让田地、户口与劳役更容易被比较和调动；它也可能把负担加在原本已受战乱、迁徙或边防牵累的家庭身上。地方中介既能使制度落地，也能借掌握信息、仓储或武力改变其样子。所谓成功，不宜只看诏令颁布或城池攻下，还要看这种接合能否持续。

本节年序主要依据《晋书》帝纪、列传、载记，《华阳国志》、十六国史辑佚与《资治通鉴》。这些材料足以确认三五七年苻坚即位，王猛参与整顿及其前后关系，使苻坚与王猛的关中进入可核对的政治脊柱；但成书时间、编纂立场与保存条件并不相同。官修史书长于保存年序、政权语言与显贵行动，地方材料长于锚定特定人地；它们各有沉默，不能把诏令、兵数、族名或人物动机直接扩展成全社会的事实。 因此，叙述把能够互证的政权转换、行军、任命和城市易手写成事实，把兵数、伤亡、族属边界、地方响应和人物动机保留为需要继续检验的问题。

墓志、题记、文书、城址与道路遗存提供的是另一种尺度。它们有时能够校正姓名、日期、官职和特定的社会网络，也能让都城叙事之外的材料进入视野；它们却不能单独代表整座城市、一个州郡或所有流民。把这两类材料放在一起，不是为了制造一幅无缝全景，而是为了知道哪些判断可以成立，哪些地方仍应保留空白。

从这个角度回看，苻坚与王猛的关中不是供读者逐条勾选的事件卡，而是持续展开的历史处境。簿册、仓粟、城门和任命使国家能力可见，也使它的脆弱处可见：只要文书、军粮、人员和道路不能重新接合，皇帝名义便不能替代地方日常。后来的政权会以新制度、宗教安排、军事编制或城市工程继续回答这一难题，同时也会制造新的受益者、成本和不确定性。

\subsection{前秦吞并诸国}
这里最重要的不是替后人判断成败，而是辨认当时有哪些资源已经断裂或仍可使用。 三七〇至三七六年前秦连灭前燕前凉等。河北、陇右、河西、汉中、代北中的新年号、仓储、驿路和军队不是背景，而是命令、粮食、人员和消息能够移动的条件。苻坚、王猛、旧国官员军人、降附者被放在同一条不平坦的关系链上：有人借它发布号令，有人依它转运供给，也有人在它被截断时改变依附。征服速度快于整合速度，广阔资源也带来协调和供给负担。

若只把三七〇至三七六年前秦连灭前燕前凉等归结为一位皇帝、将领或谋臣的性格，便会错过更长的压力。继承安排、军镇分布、土地关系、边防开支和迁徙，往往早已改变了地方的承受能力；眼前的冲突只是把它们集中显露。政治语言可以称勤王、讨逆或归附，成本仍由持兵者、供给者和居民共同承担。 因而名义、武力和供给必须并置叙述，不能把其中任何一项当成万能的解释。

从地方尺度看，事情并不只发生在帝纪最常书写的宫城或战场。守令、吏员、舟人、农户、工匠、商人和随军家属未必留下姓名，却要承担征发、迁移、停市、守城或重登簿册的后果。仓储是否有粮，渡口是否可用，旧籍能否辨认，新来者与旧居民能否取得暂时位置，决定命令能否成为日常秩序。

人事转折也应在这种约束中理解。苻坚、王猛、旧国官员军人、降附者不是脱离制度而自由行动的人物，他们受官号、部众、家族、道路、既得任命以及对对手力量的有限消息所限。传记常把胜者写得果断、败者写得必然失败；它保存冲突轮廓，却容易遮住联盟为何在此刻成立，又为何在城防、补给或继承改变后瓦解。

三七〇至三七六年前秦连灭前燕前凉等的直接结果，通常不是把一片地方从混乱变为整齐。它会重写谁守备、哪些户口重登、哪条水道由谁巡护、哪些旧官仍可使用，也会让一部分人带着家属、技术、经卷或军籍离开原位。新获土地需要守军，新接军队需要粮饷，新立皇统需要文书；这些持续成本很快进入下一场危机，形成难以回避的反馈回路。

这也是制度得失必须落在具体人事上的原因。能够提高国家可见性的制度，可能让田地、户口与劳役更容易被比较和调动；它也可能把负担加在原本已受战乱、迁徙或边防牵累的家庭身上。地方中介既能使制度落地，也能借掌握信息、仓储或武力改变其样子。所谓成功，不宜只看诏令颁布或城池攻下，还要看这种接合能否持续。

本节年序主要依据《晋书》帝纪、列传、载记，《华阳国志》、十六国史辑佚与《资治通鉴》。这些材料足以确认三七〇至三七六年前秦连灭前燕前凉等及其前后关系，使前秦吞并诸国进入可核对的政治脊柱；但成书时间、编纂立场与保存条件并不相同。官修史书长于保存年序、政权语言与显贵行动，地方材料长于锚定特定人地；它们各有沉默，不能把诏令、兵数、族名或人物动机直接扩展成全社会的事实。 因此，叙述把能够互证的政权转换、行军、任命和城市易手写成事实，把兵数、伤亡、族属边界、地方响应和人物动机保留为需要继续检验的问题。

墓志、题记、文书、城址与道路遗存提供的是另一种尺度。它们有时能够校正姓名、日期、官职和特定的社会网络，也能让都城叙事之外的材料进入视野；它们却不能单独代表整座城市、一个州郡或所有流民。把这两类材料放在一起，不是为了制造一幅无缝全景，而是为了知道哪些判断可以成立，哪些地方仍应保留空白。

从这个角度回看，前秦吞并诸国不是供读者逐条勾选的事件卡，而是持续展开的历史处境。新年号、仓储、驿路和军队使国家能力可见，也使它的脆弱处可见：只要文书、军粮、人员和道路不能重新接合，皇帝名义便不能替代地方日常。后来的政权会以新制度、宗教安排、军事编制或城市工程继续回答这一难题，同时也会制造新的受益者、成本和不确定性。

\subsection{桓温北伐的限度}
现存材料不能还原每一个现场，却能让人看见选择被空间和供给怎样限定。 三五六年桓温入洛阳，三六九年枋头失利。荆州、洛阳、黄河渡口与枋头中的帝陵、粮船、营垒和退军路不是背景，而是命令、粮食、人员和消息能够移动的条件。桓温、东晋军、地方接应者、北方政权被放在同一条不平坦的关系链上：有人借它发布号令，有人依它转运供给，也有人在它被截断时改变依附。收复旧都象征不能替代长期粮道和驻军。

若只把三五六年桓温入洛阳，三六九年枋头失利归结为一位皇帝、将领或谋臣的性格，便会错过更长的压力。继承安排、军镇分布、土地关系、边防开支和迁徙，往往早已改变了地方的承受能力；眼前的冲突只是把它们集中显露。位置、资源、联盟和有限消息限定选择，未被记录的心理不应由叙述替他们补足。 因而名义、武力和供给必须并置叙述，不能把其中任何一项当成万能的解释。

从地方尺度看，事情并不只发生在帝纪最常书写的宫城或战场。守令、吏员、舟人、农户、工匠、商人和随军家属未必留下姓名，却要承担征发、迁移、停市、守城或重登簿册的后果。仓储是否有粮，渡口是否可用，旧籍能否辨认，新来者与旧居民能否取得暂时位置，决定命令能否成为日常秩序。

人事转折也应在这种约束中理解。桓温、东晋军、地方接应者、北方政权不是脱离制度而自由行动的人物，他们受官号、部众、家族、道路、既得任命以及对对手力量的有限消息所限。传记常把胜者写得果断、败者写得必然失败；它保存冲突轮廓，却容易遮住联盟为何在此刻成立，又为何在城防、补给或继承改变后瓦解。

三五六年桓温入洛阳，三六九年枋头失利的直接结果，通常不是把一片地方从混乱变为整齐。它会重写谁守备、哪些户口重登、哪条水道由谁巡护、哪些旧官仍可使用，也会让一部分人带着家属、技术、经卷或军籍离开原位。新获土地需要守军，新接军队需要粮饷，新立皇统需要文书；这些持续成本很快进入下一场危机，形成难以回避的反馈回路。

这也是制度得失必须落在具体人事上的原因。能够提高国家可见性的制度，可能让田地、户口与劳役更容易被比较和调动；它也可能把负担加在原本已受战乱、迁徙或边防牵累的家庭身上。地方中介既能使制度落地，也能借掌握信息、仓储或武力改变其样子。所谓成功，不宜只看诏令颁布或城池攻下，还要看这种接合能否持续。

本节年序主要依据《晋书》帝纪、列传、载记，《华阳国志》、十六国史辑佚与《资治通鉴》。这些材料足以确认三五六年桓温入洛阳，三六九年枋头失利及其前后关系，使桓温北伐的限度进入可核对的政治脊柱；但成书时间、编纂立场与保存条件并不相同。官修史书长于保存年序、政权语言与显贵行动，地方材料长于锚定特定人地；它们各有沉默，不能把诏令、兵数、族名或人物动机直接扩展成全社会的事实。 因此，叙述把能够互证的政权转换、行军、任命和城市易手写成事实，把兵数、伤亡、族属边界、地方响应和人物动机保留为需要继续检验的问题。

墓志、题记、文书、城址与道路遗存提供的是另一种尺度。它们有时能够校正姓名、日期、官职和特定的社会网络，也能让都城叙事之外的材料进入视野；它们却不能单独代表整座城市、一个州郡或所有流民。把这两类材料放在一起，不是为了制造一幅无缝全景，而是为了知道哪些判断可以成立，哪些地方仍应保留空白。

从这个角度回看，桓温北伐的限度不是供读者逐条勾选的事件卡，而是持续展开的历史处境。帝陵、粮船、营垒和退军路使国家能力可见，也使它的脆弱处可见：只要文书、军粮、人员和道路不能重新接合，皇帝名义便不能替代地方日常。后来的政权会以新制度、宗教安排、军事编制或城市工程继续回答这一难题，同时也会制造新的受益者、成本和不确定性。

\subsection{淝水前的两种军队}
把时间倒回当时，行动者面对的从来不是后来已经知道的结局。 前秦南伐与北府兵在淮河相遇。京口、寿阳、淮河和寿春中的军船、远征辎重、渡口和粮仓不是背景，而是命令、粮食、人员和消息能够移动的条件。谢玄、谢石、苻坚、诸将、转运者被放在同一条不平坦的关系链上：有人借它发布号令，有人依它转运供给，也有人在它被截断时改变依附。近本土水陆协同与远距离多来源补给，构成双方不同限制。

若只把前秦南伐与北府兵在淮河相遇归结为一位皇帝、将领或谋臣的性格，便会错过更长的压力。继承安排、军镇分布、土地关系、边防开支和迁徙，往往早已改变了地方的承受能力；眼前的冲突只是把它们集中显露。能署名的诏令未必有军队，掌兵者也未必得到广泛承认。 因而名义、武力和供给必须并置叙述，不能把其中任何一项当成万能的解释。

从地方尺度看，事情并不只发生在帝纪最常书写的宫城或战场。守令、吏员、舟人、农户、工匠、商人和随军家属未必留下姓名，却要承担征发、迁移、停市、守城或重登簿册的后果。仓储是否有粮，渡口是否可用，旧籍能否辨认，新来者与旧居民能否取得暂时位置，决定命令能否成为日常秩序。

人事转折也应在这种约束中理解。谢玄、谢石、苻坚、诸将、转运者不是脱离制度而自由行动的人物，他们受官号、部众、家族、道路、既得任命以及对对手力量的有限消息所限。传记常把胜者写得果断、败者写得必然失败；它保存冲突轮廓，却容易遮住联盟为何在此刻成立，又为何在城防、补给或继承改变后瓦解。

前秦南伐与北府兵在淮河相遇的直接结果，通常不是把一片地方从混乱变为整齐。它会重写谁守备、哪些户口重登、哪条水道由谁巡护、哪些旧官仍可使用，也会让一部分人带着家属、技术、经卷或军籍离开原位。新获土地需要守军，新接军队需要粮饷，新立皇统需要文书；这些持续成本很快进入下一场危机，形成难以回避的反馈回路。

这也是制度得失必须落在具体人事上的原因。能够提高国家可见性的制度，可能让田地、户口与劳役更容易被比较和调动；它也可能把负担加在原本已受战乱、迁徙或边防牵累的家庭身上。地方中介既能使制度落地，也能借掌握信息、仓储或武力改变其样子。所谓成功，不宜只看诏令颁布或城池攻下，还要看这种接合能否持续。

本节年序主要依据《晋书》帝纪、列传、载记，《华阳国志》、十六国史辑佚与《资治通鉴》。这些材料足以确认前秦南伐与北府兵在淮河相遇及其前后关系，使淝水前的两种军队进入可核对的政治脊柱；但成书时间、编纂立场与保存条件并不相同。官修史书长于保存年序、政权语言与显贵行动，地方材料长于锚定特定人地；它们各有沉默，不能把诏令、兵数、族名或人物动机直接扩展成全社会的事实。 因此，叙述把能够互证的政权转换、行军、任命和城市易手写成事实，把兵数、伤亡、族属边界、地方响应和人物动机保留为需要继续检验的问题。

墓志、题记、文书、城址与道路遗存提供的是另一种尺度。它们有时能够校正姓名、日期、官职和特定的社会网络，也能让都城叙事之外的材料进入视野；它们却不能单独代表整座城市、一个州郡或所有流民。把这两类材料放在一起，不是为了制造一幅无缝全景，而是为了知道哪些判断可以成立，哪些地方仍应保留空白。

从这个角度回看，淝水前的两种军队不是供读者逐条勾选的事件卡，而是持续展开的历史处境。军船、远征辎重、渡口和粮仓使国家能力可见，也使它的脆弱处可见：只要文书、军粮、人员和道路不能重新接合，皇帝名义便不能替代地方日常。后来的政权会以新制度、宗教安排、军事编制或城市工程继续回答这一难题，同时也会制造新的受益者、成本和不确定性。

\subsection{淝水与前秦瓦解}
从一件可见之物进入，才不会把这段历史写成抽象的王朝更替。 三八三年东晋胜于淝水，前秦随即分裂。淝水、寿阳、淮南和北方退路中的河岸阵地、辎重、战报和城门不是背景，而是命令、粮食、人员和消息能够移动的条件。谢玄、苻坚、前秦军士、东晋守军被放在同一条不平坦的关系链上：有人借它发布号令，有人依它转运供给，也有人在它被截断时改变依附。战役打破中枢协调，原被征服地区的军政网络重新获得空间。

若只把三八三年东晋胜于淝水，前秦随即分裂归结为一位皇帝、将领或谋臣的性格，便会错过更长的压力。继承安排、军镇分布、土地关系、边防开支和迁徙，往往早已改变了地方的承受能力；眼前的冲突只是把它们集中显露。有人能守城，有人能运粮，也有人只剩下离开原处的道路。 因而名义、武力和供给必须并置叙述，不能把其中任何一项当成万能的解释。

从地方尺度看，事情并不只发生在帝纪最常书写的宫城或战场。守令、吏员、舟人、农户、工匠、商人和随军家属未必留下姓名，却要承担征发、迁移、停市、守城或重登簿册的后果。仓储是否有粮，渡口是否可用，旧籍能否辨认，新来者与旧居民能否取得暂时位置，决定命令能否成为日常秩序。

人事转折也应在这种约束中理解。谢玄、苻坚、前秦军士、东晋守军不是脱离制度而自由行动的人物，他们受官号、部众、家族、道路、既得任命以及对对手力量的有限消息所限。传记常把胜者写得果断、败者写得必然失败；它保存冲突轮廓，却容易遮住联盟为何在此刻成立，又为何在城防、补给或继承改变后瓦解。

三八三年东晋胜于淝水，前秦随即分裂的直接结果，通常不是把一片地方从混乱变为整齐。它会重写谁守备、哪些户口重登、哪条水道由谁巡护、哪些旧官仍可使用，也会让一部分人带着家属、技术、经卷或军籍离开原位。新获土地需要守军，新接军队需要粮饷，新立皇统需要文书；这些持续成本很快进入下一场危机，形成难以回避的反馈回路。

这也是制度得失必须落在具体人事上的原因。能够提高国家可见性的制度，可能让田地、户口与劳役更容易被比较和调动；它也可能把负担加在原本已受战乱、迁徙或边防牵累的家庭身上。地方中介既能使制度落地，也能借掌握信息、仓储或武力改变其样子。所谓成功，不宜只看诏令颁布或城池攻下，还要看这种接合能否持续。

本节年序主要依据《晋书》帝纪、列传、载记，《华阳国志》、十六国史辑佚与《资治通鉴》。这些材料足以确认三八三年东晋胜于淝水，前秦随即分裂及其前后关系，使淝水与前秦瓦解进入可核对的政治脊柱；但成书时间、编纂立场与保存条件并不相同。官修史书长于保存年序、政权语言与显贵行动，地方材料长于锚定特定人地；它们各有沉默，不能把诏令、兵数、族名或人物动机直接扩展成全社会的事实。 因此，叙述把能够互证的政权转换、行军、任命和城市易手写成事实，把兵数、伤亡、族属边界、地方响应和人物动机保留为需要继续检验的问题。

墓志、题记、文书、城址与道路遗存提供的是另一种尺度。它们有时能够校正姓名、日期、官职和特定的社会网络，也能让都城叙事之外的材料进入视野；它们却不能单独代表整座城市、一个州郡或所有流民。把这两类材料放在一起，不是为了制造一幅无缝全景，而是为了知道哪些判断可以成立，哪些地方仍应保留空白。

从这个角度回看，淝水与前秦瓦解不是供读者逐条勾选的事件卡，而是持续展开的历史处境。河岸阵地、辎重、战报和城门使国家能力可见，也使它的脆弱处可见：只要文书、军粮、人员和道路不能重新接合，皇帝名义便不能替代地方日常。后来的政权会以新制度、宗教安排、军事编制或城市工程继续回答这一难题，同时也会制造新的受益者、成本和不确定性。

\subsection{刘裕的关中短留}
这里最重要的不是替后人判断成败，而是辨认当时有哪些资源已经断裂或仍可使用。 四一六至四一八年刘裕入洛阳长安，关中旋失。江淮、洛阳、潼关、长安和关陇中的船队、潼关、仓储和驻军不是背景，而是命令、粮食、人员和消息能够移动的条件。刘裕、后秦旧部、赫连勃勃、地方力量被放在同一条不平坦的关系链上：有人借它发布号令，有人依它转运供给，也有人在它被截断时改变依附。南方仍能远征，前线久守受粮运、接应和军镇利益限制。

若只把四一六至四一八年刘裕入洛阳长安，关中旋失归结为一位皇帝、将领或谋臣的性格，便会错过更长的压力。继承安排、军镇分布、土地关系、边防开支和迁徙，往往早已改变了地方的承受能力；眼前的冲突只是把它们集中显露。政治语言可以称勤王、讨逆或归附，成本仍由持兵者、供给者和居民共同承担。 因而名义、武力和供给必须并置叙述，不能把其中任何一项当成万能的解释。

从地方尺度看，事情并不只发生在帝纪最常书写的宫城或战场。守令、吏员、舟人、农户、工匠、商人和随军家属未必留下姓名，却要承担征发、迁移、停市、守城或重登簿册的后果。仓储是否有粮，渡口是否可用，旧籍能否辨认，新来者与旧居民能否取得暂时位置，决定命令能否成为日常秩序。

人事转折也应在这种约束中理解。刘裕、后秦旧部、赫连勃勃、地方力量不是脱离制度而自由行动的人物，他们受官号、部众、家族、道路、既得任命以及对对手力量的有限消息所限。传记常把胜者写得果断、败者写得必然失败；它保存冲突轮廓，却容易遮住联盟为何在此刻成立，又为何在城防、补给或继承改变后瓦解。

四一六至四一八年刘裕入洛阳长安，关中旋失的直接结果，通常不是把一片地方从混乱变为整齐。它会重写谁守备、哪些户口重登、哪条水道由谁巡护、哪些旧官仍可使用，也会让一部分人带着家属、技术、经卷或军籍离开原位。新获土地需要守军，新接军队需要粮饷，新立皇统需要文书；这些持续成本很快进入下一场危机，形成难以回避的反馈回路。

这也是制度得失必须落在具体人事上的原因。能够提高国家可见性的制度，可能让田地、户口与劳役更容易被比较和调动；它也可能把负担加在原本已受战乱、迁徙或边防牵累的家庭身上。地方中介既能使制度落地，也能借掌握信息、仓储或武力改变其样子。所谓成功，不宜只看诏令颁布或城池攻下，还要看这种接合能否持续。

本节年序主要依据《晋书》帝纪、列传、载记，《华阳国志》、十六国史辑佚与《资治通鉴》。这些材料足以确认四一六至四一八年刘裕入洛阳长安，关中旋失及其前后关系，使刘裕的关中短留进入可核对的政治脊柱；但成书时间、编纂立场与保存条件并不相同。官修史书长于保存年序、政权语言与显贵行动，地方材料长于锚定特定人地；它们各有沉默，不能把诏令、兵数、族名或人物动机直接扩展成全社会的事实。 因此，叙述把能够互证的政权转换、行军、任命和城市易手写成事实，把兵数、伤亡、族属边界、地方响应和人物动机保留为需要继续检验的问题。

墓志、题记、文书、城址与道路遗存提供的是另一种尺度。它们有时能够校正姓名、日期、官职和特定的社会网络，也能让都城叙事之外的材料进入视野；它们却不能单独代表整座城市、一个州郡或所有流民。把这两类材料放在一起，不是为了制造一幅无缝全景，而是为了知道哪些判断可以成立，哪些地方仍应保留空白。

从这个角度回看，刘裕的关中短留不是供读者逐条勾选的事件卡，而是持续展开的历史处境。船队、潼关、仓储和驻军使国家能力可见，也使它的脆弱处可见：只要文书、军粮、人员和道路不能重新接合，皇帝名义便不能替代地方日常。后来的政权会以新制度、宗教安排、军事编制或城市工程继续回答这一难题，同时也会制造新的受益者、成本和不确定性。

\begin{causalchain}[本节制度得失]
这些事件显示，最高政权的转换只能提供重新组织的起点。能否把官员、户籍、军队、仓储、宗教机构与地方社会重新接合，才决定新秩序能维持多久；这种接合的成本和不均衡，也会成为下一次危机的条件。
\end{causalchain}

% 深描续写：制度、社会与材料的连续尺度
\section{大规模征服为何难以变成稳定控制}
第一，时间上的连续并不等于政治上的连续。一个年号被替换，或一座城的城门被打开，并不自动让原有的债务、亲族、服役关系和地方记忆消失；它们常在新政权的名目下继续运作，也会在新的征发中显出抵触。 在关中、河北、淮河、洛阳、长安和代北这一组相互连通又彼此牵制的地点，军船、屯田、河岸、辎重、城戍和潼关把抽象的权力落到可以触摸的层面。前秦扩张、桓温北伐、淝水、后燕后秦北魏兴起和刘裕北伐并非孤立的名词串，而是让苻坚、王猛、谢玄、桓温、慕容垂、姚苌、刘裕不得不在资源有限、消息不全的条件下重新排列位置的过程。

就大规模征服为何难以变成稳定控制而言，关键不在于把每个变化归还给一个唯一原因，而在于追问它怎样与前后事件相接。南北朝由胜利之后仍无法接收地方的困境展开 这种前后相承不是机械决定论：同样的制度或道路，在不同的都城、边境和乡里可能产生不同结果，因为掌握中介、承担成本和接收消息的人并不相同。

年序与主要行动可由《晋书》帝纪列传载记、《华阳国志》、十六国史辑佚与《资治通鉴》相互校读；但史书中的兵数、死亡、动机、族名和地方响应经常带有修辞或编纂立场。墓志、题记、文书、城址和交通遗存可以使具体姓名、日期与物质网络更清楚，却同样不能代替社会全景。因而这一节只在材料足以支撑之处推进解释，在无法互证之处保留疑问。

第二，空间必须被当作历史行动者。河流、关隘、山道、港口和城市并不只规定军队能否通过，也规定消息怎样迟到、粮食怎样损耗、逃亡者在哪里停留，以及中央命令在何处必须转交给地方中介。 在关中、河北、淮河、洛阳、长安和代北这一组相互连通又彼此牵制的地点，军船、屯田、河岸、辎重、城戍和潼关把抽象的权力落到可以触摸的层面。前秦扩张、桓温北伐、淝水、后燕后秦北魏兴起和刘裕北伐并非孤立的名词串，而是让苻坚、王猛、谢玄、桓温、慕容垂、姚苌、刘裕不得不在资源有限、消息不全的条件下重新排列位置的过程。

就大规模征服为何难以变成稳定控制而言，关键不在于把每个变化归还给一个唯一原因，而在于追问它怎样与前后事件相接。南北朝由胜利之后仍无法接收地方的困境展开 这种前后相承不是机械决定论：同样的制度或道路，在不同的都城、边境和乡里可能产生不同结果，因为掌握中介、承担成本和接收消息的人并不相同。

年序与主要行动可由《晋书》帝纪列传载记、《华阳国志》、十六国史辑佚与《资治通鉴》相互校读；但史书中的兵数、死亡、动机、族名和地方响应经常带有修辞或编纂立场。墓志、题记、文书、城址和交通遗存可以使具体姓名、日期与物质网络更清楚，却同样不能代替社会全景。因而这一节只在材料足以支撑之处推进解释，在无法互证之处保留疑问。

第三，行政的力量来自可重复的登记，而非一次壮观的宣示。官府要把人口、土地、军役和税粮写入可以核验的簿册，便须依靠识字的吏员、熟悉乡里的中介和仍能维持的仓储道路；这些条件缺一，制度语言就会在地方变形。 在关中、河北、淮河、洛阳、长安和代北这一组相互连通又彼此牵制的地点，军船、屯田、河岸、辎重、城戍和潼关把抽象的权力落到可以触摸的层面。前秦扩张、桓温北伐、淝水、后燕后秦北魏兴起和刘裕北伐并非孤立的名词串，而是让苻坚、王猛、谢玄、桓温、慕容垂、姚苌、刘裕不得不在资源有限、消息不全的条件下重新排列位置的过程。

就大规模征服为何难以变成稳定控制而言，关键不在于把每个变化归还给一个唯一原因，而在于追问它怎样与前后事件相接。南北朝由胜利之后仍无法接收地方的困境展开 这种前后相承不是机械决定论：同样的制度或道路，在不同的都城、边境和乡里可能产生不同结果，因为掌握中介、承担成本和接收消息的人并不相同。

年序与主要行动可由《晋书》帝纪列传载记、《华阳国志》、十六国史辑佚与《资治通鉴》相互校读；但史书中的兵数、死亡、动机、族名和地方响应经常带有修辞或编纂立场。墓志、题记、文书、城址和交通遗存可以使具体姓名、日期与物质网络更清楚，却同样不能代替社会全景。因而这一节只在材料足以支撑之处推进解释，在无法互证之处保留疑问。

第四，战争中的普通人并非只作为数字出现。即使材料很少记录他们的名字，征发、避难、停市、重新登记、失去耕作时间和寻找新的保护者，仍然是政治变动得以发生而又付出代价的方式。叙事不能替他们虚构心声，却应当保留这种成本。 在关中、河北、淮河、洛阳、长安和代北这一组相互连通又彼此牵制的地点，军船、屯田、河岸、辎重、城戍和潼关把抽象的权力落到可以触摸的层面。前秦扩张、桓温北伐、淝水、后燕后秦北魏兴起和刘裕北伐并非孤立的名词串，而是让苻坚、王猛、谢玄、桓温、慕容垂、姚苌、刘裕不得不在资源有限、消息不全的条件下重新排列位置的过程。

就大规模征服为何难以变成稳定控制而言，关键不在于把每个变化归还给一个唯一原因，而在于追问它怎样与前后事件相接。南北朝由胜利之后仍无法接收地方的困境展开 这种前后相承不是机械决定论：同样的制度或道路，在不同的都城、边境和乡里可能产生不同结果，因为掌握中介、承担成本和接收消息的人并不相同。

年序与主要行动可由《晋书》帝纪列传载记、《华阳国志》、十六国史辑佚与《资治通鉴》相互校读；但史书中的兵数、死亡、动机、族名和地方响应经常带有修辞或编纂立场。墓志、题记、文书、城址和交通遗存可以使具体姓名、日期与物质网络更清楚，却同样不能代替社会全景。因而这一节只在材料足以支撑之处推进解释，在无法互证之处保留疑问。

第五，精英集团也不是可以用单一标签概括的整体。宗室、后族、侨姓、旧族、军将、僧团和地方豪强都在不同地点拥有不一样的资源；他们可以在一场危机中结盟，下一场危机又因任命、土地或供给而分裂。 在关中、河北、淮河、洛阳、长安和代北这一组相互连通又彼此牵制的地点，军船、屯田、河岸、辎重、城戍和潼关把抽象的权力落到可以触摸的层面。前秦扩张、桓温北伐、淝水、后燕后秦北魏兴起和刘裕北伐并非孤立的名词串，而是让苻坚、王猛、谢玄、桓温、慕容垂、姚苌、刘裕不得不在资源有限、消息不全的条件下重新排列位置的过程。

就大规模征服为何难以变成稳定控制而言，关键不在于把每个变化归还给一个唯一原因，而在于追问它怎样与前后事件相接。南北朝由胜利之后仍无法接收地方的困境展开 这种前后相承不是机械决定论：同样的制度或道路，在不同的都城、边境和乡里可能产生不同结果，因为掌握中介、承担成本和接收消息的人并不相同。

年序与主要行动可由《晋书》帝纪列传载记、《华阳国志》、十六国史辑佚与《资治通鉴》相互校读；但史书中的兵数、死亡、动机、族名和地方响应经常带有修辞或编纂立场。墓志、题记、文书、城址和交通遗存可以使具体姓名、日期与物质网络更清楚，却同样不能代替社会全景。因而这一节只在材料足以支撑之处推进解释，在无法互证之处保留疑问。

第六，军事胜利的后果总要经过安置。俘虏、降军、旧吏、流民和被征服地区的仓储，并不会因为战报写下胜利就自动变为新王朝的力量；如何给粮、编伍、授官、迁置或准其留居，往往决定征服能否持续。 在关中、河北、淮河、洛阳、长安和代北这一组相互连通又彼此牵制的地点，军船、屯田、河岸、辎重、城戍和潼关把抽象的权力落到可以触摸的层面。前秦扩张、桓温北伐、淝水、后燕后秦北魏兴起和刘裕北伐并非孤立的名词串，而是让苻坚、王猛、谢玄、桓温、慕容垂、姚苌、刘裕不得不在资源有限、消息不全的条件下重新排列位置的过程。

就大规模征服为何难以变成稳定控制而言，关键不在于把每个变化归还给一个唯一原因，而在于追问它怎样与前后事件相接。南北朝由胜利之后仍无法接收地方的困境展开 这种前后相承不是机械决定论：同样的制度或道路，在不同的都城、边境和乡里可能产生不同结果，因为掌握中介、承担成本和接收消息的人并不相同。

年序与主要行动可由《晋书》帝纪列传载记、《华阳国志》、十六国史辑佚与《资治通鉴》相互校读；但史书中的兵数、死亡、动机、族名和地方响应经常带有修辞或编纂立场。墓志、题记、文书、城址和交通遗存可以使具体姓名、日期与物质网络更清楚，却同样不能代替社会全景。因而这一节只在材料足以支撑之处推进解释，在无法互证之处保留疑问。

第七，财政不只是一项抽象收入。谷帛、田租、劳役、船只、马匹、器械和工匠的移动，决定一个都城或军镇能把多少命令变成现实；当这些资源被宫廷争夺或远征吸走时，边地和乡里会先感到秩序变薄。 在关中、河北、淮河、洛阳、长安和代北这一组相互连通又彼此牵制的地点，军船、屯田、河岸、辎重、城戍和潼关把抽象的权力落到可以触摸的层面。前秦扩张、桓温北伐、淝水、后燕后秦北魏兴起和刘裕北伐并非孤立的名词串，而是让苻坚、王猛、谢玄、桓温、慕容垂、姚苌、刘裕不得不在资源有限、消息不全的条件下重新排列位置的过程。

就大规模征服为何难以变成稳定控制而言，关键不在于把每个变化归还给一个唯一原因，而在于追问它怎样与前后事件相接。南北朝由胜利之后仍无法接收地方的困境展开 这种前后相承不是机械决定论：同样的制度或道路，在不同的都城、边境和乡里可能产生不同结果，因为掌握中介、承担成本和接收消息的人并不相同。

年序与主要行动可由《晋书》帝纪列传载记、《华阳国志》、十六国史辑佚与《资治通鉴》相互校读；但史书中的兵数、死亡、动机、族名和地方响应经常带有修辞或编纂立场。墓志、题记、文书、城址和交通遗存可以使具体姓名、日期与物质网络更清楚，却同样不能代替社会全景。因而这一节只在材料足以支撑之处推进解释，在无法互证之处保留疑问。

第八，宗教材料让国家之外的组织变得可见。寺院、道观、译场、石窟和造像把施主、工匠、经卷、田产和交通连在一起；王室的赞助或禁断当然重要，但不能替代对地方信众、财物和营造节奏的具体追问。 在关中、河北、淮河、洛阳、长安和代北这一组相互连通又彼此牵制的地点，军船、屯田、河岸、辎重、城戍和潼关把抽象的权力落到可以触摸的层面。前秦扩张、桓温北伐、淝水、后燕后秦北魏兴起和刘裕北伐并非孤立的名词串，而是让苻坚、王猛、谢玄、桓温、慕容垂、姚苌、刘裕不得不在资源有限、消息不全的条件下重新排列位置的过程。

就大规模征服为何难以变成稳定控制而言，关键不在于把每个变化归还给一个唯一原因，而在于追问它怎样与前后事件相接。南北朝由胜利之后仍无法接收地方的困境展开 这种前后相承不是机械决定论：同样的制度或道路，在不同的都城、边境和乡里可能产生不同结果，因为掌握中介、承担成本和接收消息的人并不相同。

年序与主要行动可由《晋书》帝纪列传载记、《华阳国志》、十六国史辑佚与《资治通鉴》相互校读；但史书中的兵数、死亡、动机、族名和地方响应经常带有修辞或编纂立场。墓志、题记、文书、城址和交通遗存可以使具体姓名、日期与物质网络更清楚，却同样不能代替社会全景。因而这一节只在材料足以支撑之处推进解释，在无法互证之处保留疑问。

第九，族称与地域称谓首先是当时军政和书写的分类。它们可以帮助辨认征发、赏罚、迁徙或外交发生的语境，却不应被直接转换成内部同质、永恒不变的现代群体；不同材料中的同一个词，也可能具有不同边界。 在关中、河北、淮河、洛阳、长安和代北这一组相互连通又彼此牵制的地点，军船、屯田、河岸、辎重、城戍和潼关把抽象的权力落到可以触摸的层面。前秦扩张、桓温北伐、淝水、后燕后秦北魏兴起和刘裕北伐并非孤立的名词串，而是让苻坚、王猛、谢玄、桓温、慕容垂、姚苌、刘裕不得不在资源有限、消息不全的条件下重新排列位置的过程。

就大规模征服为何难以变成稳定控制而言，关键不在于把每个变化归还给一个唯一原因，而在于追问它怎样与前后事件相接。南北朝由胜利之后仍无法接收地方的困境展开 这种前后相承不是机械决定论：同样的制度或道路，在不同的都城、边境和乡里可能产生不同结果，因为掌握中介、承担成本和接收消息的人并不相同。

年序与主要行动可由《晋书》帝纪列传载记、《华阳国志》、十六国史辑佚与《资治通鉴》相互校读；但史书中的兵数、死亡、动机、族名和地方响应经常带有修辞或编纂立场。墓志、题记、文书、城址和交通遗存可以使具体姓名、日期与物质网络更清楚，却同样不能代替社会全景。因而这一节只在材料足以支撑之处推进解释，在无法互证之处保留疑问。

第十，因果关系必须保留层次。结构性压力可能长期存在，近期触发也许只是一次废立、叛乱或失守；使行动成为可能的往往是道路、仓储、联盟和信息。把这些层次分开，才能避免把时间上相邻的变化误写成唯一原因。 在关中、河北、淮河、洛阳、长安和代北这一组相互连通又彼此牵制的地点，军船、屯田、河岸、辎重、城戍和潼关把抽象的权力落到可以触摸的层面。前秦扩张、桓温北伐、淝水、后燕后秦北魏兴起和刘裕北伐并非孤立的名词串，而是让苻坚、王猛、谢玄、桓温、慕容垂、姚苌、刘裕不得不在资源有限、消息不全的条件下重新排列位置的过程。

就大规模征服为何难以变成稳定控制而言，关键不在于把每个变化归还给一个唯一原因，而在于追问它怎样与前后事件相接。南北朝由胜利之后仍无法接收地方的困境展开 这种前后相承不是机械决定论：同样的制度或道路，在不同的都城、边境和乡里可能产生不同结果，因为掌握中介、承担成本和接收消息的人并不相同。

年序与主要行动可由《晋书》帝纪列传载记、《华阳国志》、十六国史辑佚与《资治通鉴》相互校读；但史书中的兵数、死亡、动机、族名和地方响应经常带有修辞或编纂立场。墓志、题记、文书、城址和交通遗存可以使具体姓名、日期与物质网络更清楚，却同样不能代替社会全景。因而这一节只在材料足以支撑之处推进解释，在无法互证之处保留疑问。

第十一，人物的选择值得叙述，却不应被神化。史料能让我们看见他们起兵、入城、纳降、逃奔或颁令，也能观察行动造成的后果；对未记录的密谈和心思，最诚实的写法是说明证据到此为止。 在关中、河北、淮河、洛阳、长安和代北这一组相互连通又彼此牵制的地点，军船、屯田、河岸、辎重、城戍和潼关把抽象的权力落到可以触摸的层面。前秦扩张、桓温北伐、淝水、后燕后秦北魏兴起和刘裕北伐并非孤立的名词串，而是让苻坚、王猛、谢玄、桓温、慕容垂、姚苌、刘裕不得不在资源有限、消息不全的条件下重新排列位置的过程。

就大规模征服为何难以变成稳定控制而言，关键不在于把每个变化归还给一个唯一原因，而在于追问它怎样与前后事件相接。南北朝由胜利之后仍无法接收地方的困境展开 这种前后相承不是机械决定论：同样的制度或道路，在不同的都城、边境和乡里可能产生不同结果，因为掌握中介、承担成本和接收消息的人并不相同。

年序与主要行动可由《晋书》帝纪列传载记、《华阳国志》、十六国史辑佚与《资治通鉴》相互校读；但史书中的兵数、死亡、动机、族名和地方响应经常带有修辞或编纂立场。墓志、题记、文书、城址和交通遗存可以使具体姓名、日期与物质网络更清楚，却同样不能代替社会全景。因而这一节只在材料足以支撑之处推进解释，在无法互证之处保留疑问。

第十二，后来的史书总在已知结局的立场上整理早期事件。它们提供不可替代的年序和政治词汇，也容易把成功者的选择写成必然。用同时代或地方材料提出反问，不是取消正史，而是避免让单一视角垄断问题。 在关中、河北、淮河、洛阳、长安和代北这一组相互连通又彼此牵制的地点，军船、屯田、河岸、辎重、城戍和潼关把抽象的权力落到可以触摸的层面。前秦扩张、桓温北伐、淝水、后燕后秦北魏兴起和刘裕北伐并非孤立的名词串，而是让苻坚、王猛、谢玄、桓温、慕容垂、姚苌、刘裕不得不在资源有限、消息不全的条件下重新排列位置的过程。

就大规模征服为何难以变成稳定控制而言，关键不在于把每个变化归还给一个唯一原因，而在于追问它怎样与前后事件相接。南北朝由胜利之后仍无法接收地方的困境展开 这种前后相承不是机械决定论：同样的制度或道路，在不同的都城、边境和乡里可能产生不同结果，因为掌握中介、承担成本和接收消息的人并不相同。

年序与主要行动可由《晋书》帝纪列传载记、《华阳国志》、十六国史辑佚与《资治通鉴》相互校读；但史书中的兵数、死亡、动机、族名和地方响应经常带有修辞或编纂立场。墓志、题记、文书、城址和交通遗存可以使具体姓名、日期与物质网络更清楚，却同样不能代替社会全景。因而这一节只在材料足以支撑之处推进解释，在无法互证之处保留疑问。

第十三，制度得失不能只用道德词评价。一个办法可能提高征发、守备或城市管理的能力，同时把负担转给某些家庭、军户或地方中介；这种副作用并不证明制度毫无作用，却解释了它为什么会在新的环境中引发反弹。 在关中、河北、淮河、洛阳、长安和代北这一组相互连通又彼此牵制的地点，军船、屯田、河岸、辎重、城戍和潼关把抽象的权力落到可以触摸的层面。前秦扩张、桓温北伐、淝水、后燕后秦北魏兴起和刘裕北伐并非孤立的名词串，而是让苻坚、王猛、谢玄、桓温、慕容垂、姚苌、刘裕不得不在资源有限、消息不全的条件下重新排列位置的过程。

就大规模征服为何难以变成稳定控制而言，关键不在于把每个变化归还给一个唯一原因，而在于追问它怎样与前后事件相接。南北朝由胜利之后仍无法接收地方的困境展开 这种前后相承不是机械决定论：同样的制度或道路，在不同的都城、边境和乡里可能产生不同结果，因为掌握中介、承担成本和接收消息的人并不相同。

年序与主要行动可由《晋书》帝纪列传载记、《华阳国志》、十六国史辑佚与《资治通鉴》相互校读；但史书中的兵数、死亡、动机、族名和地方响应经常带有修辞或编纂立场。墓志、题记、文书、城址和交通遗存可以使具体姓名、日期与物质网络更清楚，却同样不能代替社会全景。因而这一节只在材料足以支撑之处推进解释，在无法互证之处保留疑问。

\begin{sourcenote}[阅读这一连续尺度]
本节不把制度、宗教、战争和迁徙拆成彼此无关的栏目。大规模征服为何难以变成稳定控制所涉及的事件，均以正史与编年材料维持最低年序，再以地方性材料限定结论范围。来源的差异不是缺陷，而是提醒读者：国家如何书写自身、地方留下何种痕迹、后来史家怎样重排叙事，回答的是不同问题。
\end{sourcenote}
